\section{Discussion}\label{sec:discussion}

\subsection{Comparison with primary occurrence studies}

Published Earth-analog occurrence estimates are not interchangeable because
their radius, orbital coordinate, HZ definition, stellar sample, catalog, and
reported statistical functional differ. The v4 median occurrence is 0.0123
and 0.0174 planets per star under constant and zero completeness, respectively
(equivalently, 1.23 and 1.74 expected planets per 100 stars, not percentages of
stars hosting planets). These values are numerically below broad-HZ estimates
as expected for the much narrower integrated phase space.

\citet{Bryson2021} use $0.5$--$1.5\,\Rearth$ and a complete conservative HZ.
Their temperature-matched 5300--6000 K G-star values of 0.38 and 0.63 are the
primary source-fidelity anchors for this analysis; the broader 4800--6300 K
$\eta_\oplus$ values of 0.37 and 0.60 provide wider context. Neither comparison
is independent corroboration. \citet{Kunimoto2020} report an 84.1\% upper limit below 0.18
planets per G star for $0.75$--$1.5\,\Rearth$ and 0.99--1.70 au; semimajor axis,
radius, and host definitions all differ from v4. The conceptually narrower
$\zeta_{1.0}$ of \citet{Burke2015} uses planet radius and orbital period within
20\% of Earth, yielding a central value near 0.1 but an allowed 0.01--2 range
once major statistical and systematic uncertainty is considered.

Using a stripped-core-aware occurrence interpretation, \citet{Pascucci2019}
obtain broad-HZ values of roughly 5--10\% and show that extrapolation choices
can change the result by factors of several. \citet{Bergsten2022} report
$\Gamma_\oplus=15^{+6}_{-4}\%$, a local differential density per logarithmic
period--radius area rather than an integrated planets-per-star quantity for the
v4 box. Converting it would require an explicit domain and occurrence model.

No one of these primary results simultaneously matches the v4 radius,
instellation, climate intersection, host-temperature range, and catalog/model
definition. Point-estimate differences therefore do not by themselves
establish inconsistency. The common conclusion relevant here is that sparse
long-period data and occurrence-model form are material epistemic limitations.

\subsection{Local support and model-form risk}

The zero-candidate audit sharpens the limitation that version~3 described only
as sparsity. The posterior remains mathematically finite because a separable
power law is fitted over the larger source domain and then integrated over the
target. Thus the interval describes uncertainty in that fitted projection; it
does not demonstrate that the local surface is smooth, separable, or free of a
break near Earth radius and instellation.

The Model~2 point comparison is useful but insufficient. It changes one
temperature parameterization within the Bryson framework and does not test
arbitrary curvature, radius--instellation interactions, or stripped-core
structure \citep{LopezRice2018,Pascucci2019}. A more complete analysis would
fit competing occurrence surfaces to the same perturbed catalogs and perform
posterior predictive checks targeted at the empty local region.

\subsection{Host-population and transport uncertainty}\label{subsec:transport}

The JJ host count is a deterministic output of the pinned parameter realization
and isochrone library. The sub-percent radial-grid change is only a numerical
diagnostic. Posterior uncertainty in star-formation and age--metallicity
histories, IMF parameters, component normalizations, radial migration, and the
choice among PARSEC, MIST, and BaSTI is not propagated. The reported host count
is reproducible but not astrophysically known to sub-percent precision.

Equation~(\ref{eq:galactic_estimator}) transports a Kepler-calibrated model to
an older Galactic population. Conditional on $\Teff$, it assumes no additional
dependence of planet occurrence or survival on age, $[\mathrm{Fe/H}]$,
$[\alpha/\mathrm{Fe}]$, disk membership, or Galactocentric radius. This is a
null transport assumption, not an empirical result.

Current empirical evidence does not justify replacing that null with one
universal correction. For close-in planets, \citet{Dai2021} find lower
$1$--$4\,\Rearth$, $P<100$ day occurrence around high-velocity stars, but the
reported separation depends on how correlated stellar properties are
controlled. \citet{BashiZucker2022} likewise find more close-in super-Earths
around their younger, slower, metal-richer thin-disk sample than around the
older thick-disk sample, while emphasizing the small thick-disk and halo
planet samples and the covariance among age, chemistry, and kinematics. In
contrast, the direct DR25 age-binned analysis of \citet{Sayeed2025} finds no
statistically significant overall occurrence--age trend for FGK stars; only a
weak decrease appears in the low-mass, metal-rich subset. These studies probe
shorter periods and much broader radius domains than the present HZ target.
Together they foreground an unresolved transport uncertainty, not a measured
multiplicative adjustment to Equation~(\ref{eq:galactic_estimator}).

The Kepler source denominator also includes observational-quality and binary
cuts that the intrinsic JJ single-star population does not reproduce. Close
companions can suppress planet occurrence \citep{Kraus2016}, while unresolved
multiplicity alters radii and completeness. The direction and magnitude of a
correction cannot be identified from the current inputs, so multiplicity
remains unquantified rather than receiving an ad hoc factor.

\subsection{Meaning of the Galactic annulus}

The 7--9 kpc interval is an absolute geometric mask motivated by historically
highlighted Galactic-habitable-zone work \citep{Lineweaver2004}; it does not
import that study's time-dependent metallicity, supernova, terrestrial-planet,
or biological weighting. Sharply bounded Galactic habitable zones are
themselves model-dependent \citep{Prantzos2008,Gowanlock2011}. The present
result is therefore an estimate within an annulus, not the number of ExoEarth
candidates in a physically complete Galactic habitable zone.

The annulus also describes present-day position rather than formation
environment. Using an action-, age-, and chemistry-based model of APOGEE
low-$\alpha$ disk stars, \citet{Frankel2020} find that angular-momentum
redistribution is substantially stronger than radial heating. Hence a star now
inside 7--9 kpc need not have formed there, and a star born there need not
remain there. Because the present calculation does not model birth radii or
planet survival along migrated orbits, it must not be interpreted as a sharply
bounded birth-annulus census.

\subsection{What the reported values mean}

The statement $\LamEE=3.22$ million in the constant branch means that the
frozen model stack assigns that median expected number of planets to the stated
radius--instellation--present-day-HZ domain and host population. The zero branch
assigns 4.57 million under a different completeness extrapolation. These are
not observed counts, not numbers of systems with at least one selected planet,
and not counts of rocky, ocean-bearing, habitable, or inhabited worlds.

The displayed posterior intervals are statistically meaningful for the
implemented conditional model, but narrower in scope than the total scientific
uncertainty. Reproducibility of a conditional projection and empirical validity
of the projection are separate questions; v4 passes the former and explicitly
fails the local-support component of the latter.
